% !TeX root = ../../Book3.tex
\section{分布观点}
\label{b3:sec:1504}

调和与次调和的差别还可以用一个正测度表达。它记录函数偏离平均等式的来源，而不要求函数具有通常的二阶导数。

\subsection{分布意义下的 Laplace 算子}
\label{b3:subsec:150401}

对于 \(u\in L^1_{\mathrm{loc}}(D)\)，定义 \(\langle\Delta u,\varphi\rangle=\int_Du\Delta\varphi\,\mathrm{d}A\)，其中 \(\varphi\in C_c^\infty(D)\)。分布配对不加复共轭，二阶积分分部使符号为正。
\begin{proposition}[对数基本解]\label{b3:prop:15-log-fundamental}
在复平面分布意义下，\(\Delta\log|z|=2\pi\delta_0\)。
\end{proposition}
\begin{proof}
删去半径 \(\varepsilon\) 的圆盘，在余域用 Green 公式。小圆的余域外法向朝向原点，边界项为
\[
\int_0^{2\pi}\bigl(-\varepsilon\log\varepsilon\,\partial_r\varphi(\varepsilon e^{\mathrm{i}t})+\varphi(\varepsilon e^{\mathrm{i}t})\bigr)\,\mathrm{d}t.
\]
第一项趋零，第二项趋 \(2\pi\varphi(0)\)。被删圆盘上 \(\log|z|\) 的可积性保证原积分极限成立。
\end{proof}

\subsection{正分布}
\label{b3:subsec:150402}

实分布 \(T\) 对每个非负测试函数取非负值时，称为\term[zheng-fenbu]{正分布}[positive distribution]。取 \(\chi\in C_c^\infty(D)\)，\(\chi=1\) 于紧集 \(K\) 邻域，则支集在 \(K\) 的实测试函数满足
\[
|T(\varphi)|\le T(\chi)\cdot\|\varphi\|_\infty.
\]
这是由 \(-\|\varphi\|_\infty\chi\le\varphi\le\|\varphi\|_\infty\chi\) 得到的。光滑逼近因此把 \(T\) 唯一延拓到连续紧支集函数，并保持正性。第一卷定理11.4.1（正泛函的 Riesz--Markov 表示）给出局部有限正 Radon 测度。这里调用的是表示定理，不把分布的高阶连续性直接误当作测度性。

\subsection{Riesz 测度}
\label{b3:subsec:150403}

\begin{proposition}\label{b3:prop:15-riesz-measure}
次调和函数 \(u\) 的分布 \(\Delta u\) 为正分布。故有唯一局部有限正测度
\(\mu_u=(2\pi)^{-1}\Delta u\)，称为其\term[riesz-cedu]{Riesz 测度}[Riesz measure]。
\end{proposition}
\begin{proof}
径向非负磨光 \(u_\varepsilon\) 在缩小的区域中仍满足平均不等式；它光滑，故 \(\Delta u_\varepsilon\ge0\)。又 \(u_\varepsilon\to u\) 于局部 \(L^1\)，对非负测试函数取极限得到 \(\int u\Delta\varphi\ge0\)。上一小节把该正分布识别为测度，唯一性由测试函数稠密性得到。
\end{proof}
例如 \(u=|z|^2\) 的测度为 \((2/\pi)\mathrm{d}A\)，而 \(u=\log|z-a|\) 的测度为 \(\delta_a\)。连续分布与点质量都可以出现，不能把 Riesz 测度只理解为离散零点。

\subsection{势论解释}
\label{b3:subsec:150404}

对光滑紧支集非负密度 \(\rho\)，对数势
\[
V(z)=\int_{\mathbb C}\log|z-w|\cdot\rho(w)\,\mathrm{d}A(w)
\]
局部可积，且 \(\Delta V=2\pi\rho\)。换序由紧支集上的对数核局部可积保证。因而若 \(\mu_u=\rho\mathrm{d}A\)，\(u-V\) 的分布 Laplace 算子为零。

分布调和函数的光滑性可以这样恢复：其磨光 \(h_\varepsilon\) 调和；在一个内部圆盘上再用固定径向核 \(\eta\)，平均值性质给 \(h_\varepsilon=h_\varepsilon*\eta\)。局部 \(L^1\) 极限使右端连同每个阶导数一致收敛，得到光滑调和代表 \(h=h*\eta\)。这是 \term[weyl-yinli]{Weyl 引理}[Weyl's lemma]在平面上的证明。

势决定非调和部分，调和函数决定不被测度看见的自由度。在一般 Radon 密度情形，对数势可能在薄集上为负无穷，必须另行选择次调和代表；这里不把局部可积函数的任意零测修改说成同一个点态次调和函数。


正分布的零阶估计还说明其局部有限性。取 \(K\Subset D\)，所有 \(0\le\varphi\le1\)、支集在 \(K\) 的测试函数均满足 \(T(\varphi)\le T(\chi)\)。表示测度在 \(K\) 的质量因此不超过同一有限常数。不同紧集选择不同 \(\chi\)；不能从局部有限性推出全域测度有限。

\ProseParagraphBreak

例如 \(u(z)=|z|^2\) 的 Riesz 测度在每个有界紧集有限，却在全平面质量无穷。\(u=\log|z|\) 的测度总质量一，但是函数在无穷远不有界。测度有限性与函数有界性是不同性质，后面 Green 势的边界控制不能只根据其中一项猜出。

\ProseParagraphBreak

正分布一定是测度，任意分布却不一定如此。点质量的导数 \(\partial_x\delta_0\) 对测试函数的作用是 \(-\partial_x\varphi(0)\)。可取支集缩小、最大模不超过一而原点导数任意大的光滑函数，所以它不满足紧集上的零阶最大范数界，不能由局部有限测度表示。它也不正：非负测试函数在原点的一阶导数可有不同符号，只要原点不是极值点。正性在测度识别中承担的正是把所有阶连续性降为零阶控制的作用。

\ProseParagraphBreak

对于无零点全纯函数，\(\mu_{\log|f|}=0\)，但函数的增长可以很大。例如 \(e^{z^3}\) 的对数模为 \(\operatorname{Re}z^3\)，调和且测度零。因此 Laplace 测度只记录势的源，并不记录全部调和自由度。要确定一个函数，还须给边界数据、增长条件或某种归一化；这也是下一章边值问题与再下一章 Green 最小性各自出现的原因。


\ProseParagraphBreak

对一般有限正测度也可以局部构造对数势。设 \(\mu\) 支集于紧集 \(K\)，定义 \(V(z)=\int_K\log|z-w|\,\mathrm{d}\mu(w)\)，允许负无穷。对任意有界集 \(L\)，
\[
\sup_{w\in K}\int_L|\log|z-w||\,\mathrm{d}A(z)<\infty,
\]
因为单位距离以内的负部由 \(2\pi\int_0^1r|\log r|\,\mathrm{d}r\) 控制，外面的正部由 \(K,L\) 的有界性控制。Tonelli 因而给 \(V\in L^1_{\mathrm{loc}}\)，尤其它不能在整个开集恒为负无穷。分布换序及对数基本解得到 \(\Delta V=2\pi\mu\)。

\ProseParagraphBreak

这个积分本身提供了适当的点态代表。对局部 \(z\) 变量，核有共同有限上界；减去该上界后用 Fatou 引理，得到 \(\limsup V(z_j)\le V(z)\)，故上半连续。每个核的圆周平均不小于中心值，对正测度积分后仍保持该不等式；若负部的换序有疑问，可先以光滑对数核逼近，再用下降极限处理。因此 \(V\) 次调和，而不只是一个 Laplace 正的分布等价类。

\ProseParagraphBreak

给定次调和 \(u\)，在两个同心内部圆盘 \(B'\Subset B\Subset D\) 上取 \(\mu=\mu_u|_{\overline B}\)，上述势满足 \(\Delta(u-V)=0\) 于 \(B\)，Weyl 引理给调和函数 \(h\)，于是 \(u=V+h\) 几乎处处。要恢复处处等式，可利用次调和函数的面积平均在半径趋零时回到自身：平均不等式给下界，上半连续性给上界，所以平均极限等于点值，包括值为负无穷的情形。双方几乎处处相等使所有面积平均相等，最终得到 \(u=V+h\) 处处成立。这是局部 Riesz 分解；把全部测度放入势时，远处边界与增长还须另外控制，本证明只承诺内部局部范围。
